% Part:sets-functions-relations
% Chapter: size-of-sets
% Section: comparing-sizes

\documentclass[../../../include/open-logic-section]{subfiles}

\begin{document}

\olfileid[tr]{sfr}{siz}{car}

\olsection{Farklı Büyüklükte Kümeler ve Cantor Teoremi}

\begin{explain}
İki kümenin aynı büyüklükte olduğu fikrini kesin bir biçimde ifade ettik.
Bir kümenin diğerinden daha küçük olduğu fikrini de kesin olarak ifade
edebiliriz. ``Daha küçük (ya da eş güçlü)'' tanımımız, kümeler arasında bir
!!{bijection} yerine ilk kümeden ikinciye bir !!{injection} gerektirecektir.
Böyle bir fonksiyon varsa ilk kümenin büyüklüğü ikinci kümeninkinden küçük ya
da ona eşittir. Sezgisel olarak bir kümeden diğerine giden bir !!{injection},
tanım kümesinin iki ayrı !!{element}ı görüntü kümesinin aynı !!{element}ına
eşlenmediği için fonksiyonun görüntü kümesinin en az tanım kümesi kadar
!!{element}a sahip olduğunu güvence altına alır.
\end{explain}

\begin{defn}
$A$, ancak ve ancak bir !!{injection} $f\colon A\to B$ varsa $B$'den
\emph{daha büyük değildir}; bunu $\cardle{A}{B}$ ile yazarız.
\end{defn}

Bunun yansımalı ve geçişli, fakat simetrik olmayan bir bağıntı olduğu açıktır
(bu bir alıştırma olarak bırakılmıştır). Ayrıca bir kümenin diğerinden
(kesin olarak) daha küçük olduğunu belirten bir kavram da tanımlayabiliriz.

\begin{defn}
$A$, ancak ve ancak bir !!{injection} $f\colon A\to B$ bulunmasına karşın
hiçbir !!{bijection} $g\colon A\to B$ yoksa $B$'den \emph{daha küçüktür};
bunu $\cardless{A}{B}$ ile yazarız. Başka bir deyişle
$\cardle{A}{B}$ ve $\cardneq{A}{B}$.
\end{defn}

Bu bağıntının yansımasız ve geçişli olduğu açıktır. (Bu bir alıştırma olarak
bırakılmıştır.) Bu gösterimle $A$ kümesinin !!{enumerable} olmasını
$\cardle{A}{\Nat}$ ile ifade edebiliriz. $A$'nın !!{nonenumerable} olmasıyla
$\cardless{\Nat}{A}$ eşdeğerliği ise kümelerin ilgili büyüklüklerinin
karşılaştırılabilirliği ve kullanılan küme teorisindeki seçim ilkeleri
altında okunmalıdır. Bu gösterim,
\oliflabeldef{sfr:siz:nen-alt:thm:nonenum-pownat}{%
\olref[sfr][siz][nen-alt]{thm:nonenum-pownat}'ı
$\cardless{\Nat}{\Pow{\Nat}}$ gözlemi olarak}{%
\olref[sfr][siz][nen]{thm:nonenum-pownat}'ı
$\cardless{\PosInt}{\Pow{\PosInt}}$ gözlemi olarak} yeniden ifade etmemizi
sağlar. Gerçekte \citet{Cantor1892}, bu son noktanın \emph{tümüyle genel}
olduğunu ispatlamıştır:

\begin{thm}[Cantor]\ollabel{thm:cantor}
Her $A$ kümesi için $\cardless{A}{\Pow{A}}$.
\end{thm}

\begin{proof}
$x\neq y$ ise kaplamsallık gereği $\{x\}\neq\{y\}$ ve dolayısıyla
$f(x)\neq f(y)$ olduğundan, $f(x)=\{x\}$ eşlemesi bir !!{injection}
$f\colon A\to\Pow{A}$'dır. Öyleyse $\cardle{A}{\Pow{A}}$.

\begin{editorial}
\olref[nen]{sec} varsa ayrıntılı ispatı; yoksa
\olref[nen-alt]{sec}'e koşut daha kısa ispatı sunuyoruz.
\end{editorial}

\oliflabeldef{sfr:siz:nen:sec}{%
  Şimdi bir !!{surjective} $g\colon A\to\Pow{A}$ fonksiyonunun, hele de
  !!{bijective} bir fonksiyonun var olamayacağını ve dolayısıyla
  $\cardneq{A}{\Pow{A}}$ olduğunu göstereceğiz. Bir
  $g\colon A\to\Pow{A}$ verildiğini varsayın. $g$ tümel olduğundan her
  $x\in A$, $g(x)\subseteq A$ alt kümesine eşlenir. $g$'nin örten
  olamayacağını gösterebiliriz. Bunun için tanımı gereği $g$'nin görüntü
  kümesinde bulunamayacak bir $\overline{A}\subseteq A$ alt kümesi tanımlayın:
  \[
  \overline{A}=\Setabs{x\in A}{x\notin g(x)}.
  \]
  $g(x)$ bütün $x\in A$ için tanımlı olduğundan $\overline{A}$ açıkça iyi
  tanımlı bir $A$ alt kümesidir. Ancak $g$'nin görüntü kümesinde bulunamaz.
  $x\in A$ gelişigüzel olsun; $\overline{A}\neq g(x)$ olduğunu göstereceğiz.
  $x\in g(x)$ ise $x\notin g(x)$ koşulunu sağlamaz ve $\overline{A}$'nın
  tanımı gereği $x\notin\overline{A}$ olur. $x\in\overline{A}$ ise
  $\overline{A}$'nın tanımlayıcı özelliğini sağlamalıdır; yani $x\in A$ ve
  $x\notin g(x)$ olur. $x$ gelişigüzel olduğundan her $x\in A$ için
  $x\in g(x)$ ancak ve ancak $x\notin\overline{A}$ olur ve dolayısıyla
  $g(x)\neq\overline{A}$. Başka bir deyişle $\overline{A}$, $g$'nin
  görüntü kümesinde bulunamaz; bu da $g$'nin örten olduğu varsayımıyla
  çelişir.}{Geriye $\cardneq{A}{\Pow{A}}$ olduğunu göstermek kalır. Tersini
  varsayarak $\cardeq{A}{\Pow{A}}$ olduğunu, yani bir !!{bijection}
  $g\colon A\to\Pow{A}$ bulunduğunu varsayın. Şimdi
  \[
    D=\Setabs{x\in A}{x\notin g(x)}
  \]
  kümesini düşünün. $D\subseteq A$, dolayısıyla $D\in\Pow{A}$ olduğuna
  dikkat edin. $g$ bir !!{bijection} olduğundan $g(y)=D$ olacak biçimde bir
  $y\in A$ vardır. Fakat şimdi
  \[
    y\in g(y) \text{ ancak ve ancak } y\in D
    \text{ ancak ve ancak } y\notin g(y)
  \]
  olur. Bu bir çelişkidir; dolayısıyla $\cardneq{A}{\Pow{A}}$.}{}
\end{proof}

\begin{explain}
\oliflabeldef{sfr:siz:nen:thm:nonenum-pownat}{\olref{thm:cantor}'ın ispatını
  \olref[nen]{thm:nonenum-pownat}'ın ispatıyla karşılaştırmak öğreticidir.
  Orada $\PosInt$ alt kümelerinin herhangi bir $Z_1$, $Z_2$, \dots{}
  listesinden, listede bulunmadığı güvence altında olan bir sayı kümesi
  $\overline{Z}$ kurulabileceğini gösterdik. Bu kümenin listede bulunmaması,
  her $n\in\PosInt$ için $n\in Z_n$ ancak ve ancak
  $n\notin\overline{Z}$ olmasından kaynaklanıyordu. Böylece $Z_n$ ile
  $\overline{Z}$'den birinin !!a{element}ı olup diğerinin olmayan bir sayı
  her zaman vardır. Burada da aynı fikri izleriz; yalnız indisler artık
  $\PosInt$ yerine $A$'nın !!{element}larıdır. $\overline{A}$, her $x\in A$
  için $g(x)$'ten farklı olacak biçimde tanımlanır; çünkü $x\in g(x)$ ancak
  ve ancak $x\notin\overline{A}$'dır. Yine her zaman $g(x)$ ile
  $\overline{A}$'dan birinin !!a{element}ı olup diğerinin olmayan bir $A$
  elemanı vardır. $\overline{Z}$ nasıl $Z_1$, $Z_2$, \dots{} listesinde
  bulunamazsa $\overline{A}$ da $g$'nin görüntü kümesinde bulunamaz.}{}

\oliflabeldef{sfr:siz:nen-alt:thm:nonenum-pownat}{\olref{thm:cantor}'ın
  ispatını \olref[nen-alt]{thm:nonenum-pownat}'ın ispatıyla karşılaştırmak
  öğreticidir. Orada $\Nat$ alt kümelerinin herhangi bir $N_0$, $N_1$,
  $N_2$, \dots{} listesinden, listede bulunmadığı güvence altında olan bir
  sayı kümesi $D$ kurabildiğimizi gösterdik. Her $n\in\Nat$ için $n\in N_n$
  ancak ve ancak $n\notin D$ olduğundan bu kümenin listede bulunmadığı
  güvence altındaydı. Burada da aynı fikri izleriz; yalnız indisler artık
  $\Nat$ yerine $A$'nın !!{element}larıdır. $D$, her $x\in A$ için
  $g(x)$'ten farklı olacak biçimde tanımlanır; çünkü $x\in g(x)$ ancak ve
  ancak $x\notin D$'dir.}{}

Bu ispatı Russell Paradoksu'nun ispatı
\olref[sfr][set][rus]{thm:russells-paradox} ile karşılaştırmak da yararlıdır.
Gerçekten Cantor Teoremi, Russell'ın kendi paradoksuna esin vermiştir.
\end{explain}

\begin{prob}
  Herhangi bir $A$ kümesi için bir !!{injection} $g\colon\Pow{A}\to A$
  olamayacağını gösterin. İpucu: $g\colon\Pow{A}\to A$'nın !!{injective}
  olduğunu varsayın. $D=\Setabs{g(B)}{B\subseteq A \text{ ve }g(B)\notin B}$
  kümesini düşünün. $x=g(D)$ olsun. Bir çelişki türetmek için $g$'nin
  !!{injective} olmasını kullanın.
\end{prob}

\end{document}
